\chapter{Segitiga dan Sudut}\label{ch:g7:triangles}

Dapatkah kamu menggambar segitiga dengan sisi $10$, $3$ dan $2$?
Dengan sudut $80^\circ$, $70^\circ$ dan $50^\circ$? Bab ini menjawab
kedua pertanyaan itu dengan dua kenyataan yang paling berguna dalam
geometri bidang: ketaksamaan segitiga dan \omterm{def:g1:addition:def}{jumlah} sudut $180^\circ$ ---
ditambah kosakata sudut yang datang bersama garis berpotongan dan
\omterm{def:g6:lines:perp}{garis sejajar}.

\section{Sudut di sekitar garis yang berpotongan}

\begin{definition}[Pasangan sudut]\label{def:g7:triangles:pairs}
\begin{itemize}
  \item Dua sudut disebut \emph{berpenyiku} jika \omterm{def:g1:addition:def}{jumlah} ukurannya
        $90^\circ$, dan \emph{berpelurus} jika \omterm{def:g1:addition:def}{jumlahnya}
        $180^\circ$.
  \item Ketika dua garis berpotongan, sudut yang saling berhadapan di
        seberang titik potongnya disebut \emph{bertolak
        belakang}\index{sudut bertolak belakang}.
  \item Ketika sebuah garis memotong dua garis lain, sudut yang
        letaknya sama pada tiap perpotongan disebut
        \emph{sudut sehadap}; sudut di antara kedua garis itu pada
        sisi yang berlawanan dari garis pemotongnya disebut
        \emph{sudut berseberangan}\index{sudut berseberangan}.
\end{itemize}
\end{definition}

\begin{theorem}[Kesamaan sudut]\label{thm:g7:triangles:equalities}
\begin{enumerate}
  \item Sudut yang \omterm{def:g7:triangles:pairs}{bertolak belakang} sama besar.
  \item Jika dua \omterm{def:g6:lines:perp}{garis sejajar}, sudut sehadapnya sama besar, dan
        \omterm{def:g7:triangles:pairs}{sudut berseberangannya} sama besar. Sebaliknya, sudut sehadap
        (atau \omterm{def:g7:triangles:pairs}{sudut berseberangan}) yang sama besar memaksa kedua garis
        itu \omterm{def:g6:lines:perp}{sejajar}.
\end{enumerate}
\end{theorem}

\begin{proof}[Bukti untuk 1]
Kedua sudut $\widehat{1}$ dan $\widehat{2}$ pada satu sisi sebuah
garis membentuk sudut lurus: $\widehat 1 + \widehat 2 = 180^\circ$.
Sudut $\widehat 3$ yang \omterm{def:g7:triangles:pairs}{bertolak belakang} dengan $\widehat 1$ juga
memenuhi $\widehat 2 + \widehat 3 = 180^\circ$ (dengan alasan yang
sama). Jadi $\widehat 1 = 180^\circ - \widehat 2 = \widehat 3$.
Bagian 2 diterima tanpa bukti pada tingkat ini.
\end{proof}

\begin{omfigure}
\begin{tikzpicture}[scale=1.0]
  \draw[thick] (-2.0,-0.9) -- (2.0,0.9);
  \draw[thick] (-2.0,0.9) -- (2.0,-0.9);
  \draw[omDef, ->] (0.75,0) arc (0:24.2:0.75);
  \draw[omDef, ->] (-0.75,0) arc (180:204.2:0.75);
  \node[omDef, font=\small] at (1.15,0.28) {$\widehat 1$};
  \node[omDef, font=\small] at (-1.15,-0.28) {$\widehat 3$};
  \draw[omThm, ->] (0.55,0.25) arc (24.2:155.8:0.6);
  \node[omThm, font=\small] at (0,0.75) {$\widehat 2$};
  \begin{scope}[xshift=7.2cm]
  \draw[thick] (-2.2,-1.0) -- (2.2,-1.0);
  \draw[thick] (-2.2,0.8) -- (2.2,0.8);
  \draw[thick, omExo] (-1.6,-1.6) -- (1.6,1.5);
  \draw[omDef, ->] (-0.28,-1.0) ++(0.75,0) arc (0:44:0.75);
  \node[omDef, font=\small] at (0.85,-0.62) {$\widehat a$};
  \draw[omDef, ->] (0.63,0.8) ++(0.75,0) arc (0:44:0.75);
  \node[omDef, font=\small] at (1.76,1.18) {$\widehat b$};
  \end{scope}
\end{tikzpicture}

{\small Kiri: \omterm{def:g7:triangles:pairs}{sudut bertolak belakang} $\widehat 1 = \widehat 3$
(masing-masing berpelurus dengan $\widehat 2$). Kanan: \omterm{def:g6:lines:perp}{garis sejajar}
yang dipotong garis ketiga --- sudut sehadap $\widehat a$ dan
$\widehat b$ sama besar.}
\end{omfigure}

\section{Jumlah sudut segitiga}

\begin{theorem}[Jumlah sudut]\label{thm:g7:triangles:sum}
Pada setiap segitiga, ketiga sudutnya berjumlah $180^\circ$.
\end{theorem}

\begin{proof}
Misalkan $ABC$ sebuah segitiga. Gambarlah garis melalui $A$ yang
\omterm{def:g6:lines:perp}{sejajar} dengan $(BC)$. Pada \omterm{def:g5:solids:def}{titik sudut} $A$, tiga sudut berjajar
sepanjang \omterm{def:g6:lines:perp}{garis sejajar} itu dan membentuk sudut lurus, $180^\circ$:
yang di tengah adalah $\widehat{A}$, dan kedua yang di luar
berseberangan dengan $\widehat B$ dan $\widehat C$ (\omterm{def:g6:lines:perp}{garis sejajar}!),
jadi keduanya sama dengan $\widehat B$ dan $\widehat C$. Karena itu
$\widehat B + \widehat A + \widehat C = 180^\circ$.
\end{proof}

\begin{omfigure}
\begin{tikzpicture}[scale=1.15]
  \coordinate (B) at (0,0);
  \coordinate (C) at (4.2,0);
  \coordinate (A) at (1.6,2.0);
  \draw[thick] (B) -- (C) -- (A) -- cycle;
  \draw[thick, omExo] (-0.3,2.0) -- (3.9,2.0);
  \node[above, font=\small] at (A) {$A$};
  \node[below left, font=\small] at (B) {$B$};
  \node[below right, font=\small] at (C) {$C$};
  \draw[omDef, ->] (0.55,0) arc (0:51.3:0.55);
  \node[omDef, font=\small] at (0.95,0.32) {$\widehat B$};
  \draw[omMeth, ->] (3.68,0) arc (180:142.4:0.52);
  \node[omMeth, font=\small] at (3.25,0.3) {$\widehat C$};
  \draw[omDef, ->] (A) ++(-0.55,0) arc (180:231.3:0.55);
  \node[omDef, font=\small] at (0.85,1.75) {$\widehat B$};
  \draw[omMeth, ->] (A) ++(0.52,0) arc (0:-37.6:0.52);
  \node[omMeth, font=\small] at (2.35,1.78) {$\widehat C$};
\end{tikzpicture}

{\small Buktinya dalam satu gambar: sepanjang \omterm{def:g6:lines:perp}{garis sejajar} dengan
$(BC)$ yang melalui $A$, sudut $\widehat B$, $\widehat A$,
$\widehat C$ (\omterm{def:g7:triangles:pairs}{sudut berseberangan} dengan warna yang cocok) memenuhi
sudut lurus.}
\end{omfigure}

\begin{example}\label{ex:g7:triangles:sum}
Dua sudut sebuah segitiga berukuran $67^\circ$ dan $58^\circ$. Sudut
yang ketiga berukuran
\[
180 - (67 + 58) = 180 - 125 = 55^\circ .
\]
Keadaan khusus yang layak dihafal: setiap sudut \omterm{def:g6:shapes:triangles}{segitiga sama sisi}
adalah $180 \div 3 = 60^\circ$; pada \omterm{def:g6:shapes:triangles}{segitiga siku-siku} kedua sudut
lancipnya berpenyiku (\omterm{def:g1:addition:def}{jumlahnya} $180 - 90 = 90^\circ$); sudut alas
\omterm{def:g6:shapes:triangles}{segitiga sama kaki} sama besar, jadi masing-masing
$\frac{180 - \text{puncak}}{2}$.
\end{example}

\begin{example}[Tidak ada segitiga dengan dua sudut siku-siku]\label{ex:g7:triangles:tworight}
Segitiga dengan dua sudut $90^\circ$ sudah memakai habis $180^\circ$,
sehingga menyisakan $0^\circ$ untuk sudut ketiganya --- mustahil.
\omterm{def:g1:addition:def}{Jumlah} sudutnya melarang banyak segitiga sekaligus.
\end{example}

\section{Ketaksamaan segitiga}

\begin{theorem}[Ketaksamaan segitiga]\label{thm:g7:triangles:inequality}
Pada setiap segitiga, setiap sisinya lebih pendek daripada \omterm{def:g1:addition:def}{jumlah} dua
sisi lainnya: untuk segitiga dengan sisi $a$, $b$, $c$,
\[
a < b + c .
\]
Sebaliknya, tiga panjang yang \emph{terbesarnya} lebih kecil daripada
\omterm{def:g1:addition:def}{jumlah} dua lainnya selalu dapat dibuat menjadi segitiga. Kalau yang
terbesar \emph{sama dengan} \omterm{def:g1:addition:def}{jumlah} dua lainnya, ``segitiga'' itu
menjadi datar: ketiga \omterm{def:g5:solids:def}{titik sudutnya} segaris.
\end{theorem}
\admitted

\begin{example}\label{ex:g7:triangles:inequality}
Dapatkah sisinya $10$, $3$ dan $2$? Ujilah yang terbesar: apakah
$10 < 3 + 2 = 5$? Tidak --- segitiga seperti itu tidak ada. Jangka
memperlihatkan sebabnya: busur berjari-jari $3$ dan $2$ yang digambar
dari kedua ujung \omterm{def:g6:lines:objects}{ruas garis} sepanjang $10$ tidak pernah bertemu.

Sisi $7$, $5$ dan $4$? Yang terbesar: $7 < 5 + 4 = 9$: ya, segitiganya
ada (memeriksa sisi terbesarnya sudah cukup).
\end{example}

\begin{omfigure}
\begin{tikzpicture}[scale=0.72]
  \draw[thick] (0,0) -- (10*0.72,0);
  \fill (0,0) circle (1.5pt) node[below, font=\small] {$A$};
  \fill (7.2,0) circle (1.5pt) node[below, font=\small] {$B$};
  \draw[omDef, thick] (2.16,0) arc (0:60:2.16);
  \draw[omDef, thick] (2.16,0) arc (0:-30:2.16);
  \draw[omThm, thick] (7.2,0) ++(-1.44,0) arc (180:120:1.44);
  \draw[omThm, thick] (7.2,0) ++(-1.44,0) arc (180:210:1.44);
  \node[omDef, font=\small] at (1.1,1.6) {jari-jari $3$};
  \node[omThm, font=\small] at (6.4,1.3) {jari-jari $2$};
\end{tikzpicture}

{\small Mencoba membangun segitiga dengan sisi $10$, $3$, $2$: kedua
busur jangkanya tetap berjauhan tanpa harapan, karena $3 + 2 < 10$.}
\end{omfigure}

\begin{method}[Keberadaan dan pelukisan segitiga]\label{met:g7:triangles:construct}
Diberikan tiga panjang:
\begin{enumerate}
  \item bandingkan yang terbesar dengan \omterm{def:g1:addition:def}{jumlah} dua lainnya; kalau ia
        tidak benar-benar lebih kecil, berhentilah: segitiganya tidak
        ada;
  \item kalau tidak, lukislah seperti pada
        \cref{met:g6:shapes:construct} (\omterm{def:g6:lines:objects}{ruas garis} alas, dua busur
        jangka);
  \item sesudah pelukisan apa pun dari sudut, periksalah dalam hati
        bahwa sudut yang diberikan berjumlah kurang dari $180^\circ$
        --- dengan sudut ketiganya melengkapi tepat $180^\circ$.
\end{enumerate}
\end{method}

\section{Latihan}

\begin{exercise}[$\star$]\label{exo:g7:triangles:1}
Sebutkan penyiku (sampai $90^\circ$) dan pelurus (sampai
$180^\circ$) dari: $30^\circ$; \ $45^\circ$; \ $72^\circ$.
\end{exercise}

\begin{exercise}[$\star$]\label{exo:g7:triangles:2}
Dua garis berpotongan; salah satu dari keempat sudutnya berukuran
$35^\circ$. Sebutkan ukuran ketiga sudut lainnya, dengan alasannya
masing-masing.
\end{exercise}

\begin{exercise}[$\star$]\label{exo:g7:triangles:3}
Hitunglah sudut ketiga sebuah segitiga yang dua sudut pertamanya
berukuran:
(a) $40^\circ$ dan $60^\circ$; \
(b) $90^\circ$ dan $28^\circ$; \
(c) $75^\circ$ dan $75^\circ$.
\end{exercise}

\begin{exercise}[$\star$]\label{exo:g7:triangles:4}
Sebuah \omterm{def:g6:shapes:triangles}{segitiga sama kaki} punya sudut puncak $40^\circ$. Hitunglah
kedua sudut alasnya. \omterm{def:g6:shapes:triangles}{Segitiga sama kaki} yang lain punya sudut alas
$70^\circ$: hitunglah sudut puncaknya.
\end{exercise}

\begin{exercise}[$\star$]\label{exo:g7:triangles:5}
Tripel mana yang dapat menjadi sisi sebuah segitiga? Berilah alasan
dengan ketaksamaan segitiga (sisi terbesarnya!):
\[
(6,\ 8,\ 12); \qquad
(5,\ 5,\ 11); \qquad
(4,\ 9,\ 13); \qquad
(7,\ 7,\ 7).
\]
\end{exercise}

\begin{exercise}[$\star$]\label{exo:g7:triangles:6}
Sebuah segitiga punya sudut $x$, $2x$ dan $3x$. Carilah $x$ dan ketiga
sudutnya. Segitiga jenis apakah itu?
\end{exercise}

\begin{exercise}[$\star$]\label{exo:g7:triangles:7}
Lukislah segitiga $ABC$ dengan $BC = 6$ cm,
$\widehat{ABC} = 50^\circ$ dan $\widehat{ACB} = 60^\circ$ (busur
derajat di $B$ dan di $C$). Sebelum melukis, ramalkan ukuran
$\widehat{BAC}$, lalu periksalah pada gambarmu.
\end{exercise}

\begin{exercise}[$\star\star$]\label{exo:g7:triangles:8}
Dua garis yang \omterm{def:g6:lines:perp}{sejajar} dipotong garis ketiga, membentuk sudut $54^\circ$
pada perpotongan yang pertama (antara garis pemotongnya dan salah satu
\omterm{def:g6:lines:perp}{garis sejajarnya}). Gambarlah keadaannya lalu sebutkan ukuran kedelapan
sudut yang terbentuk.
\end{exercise}

\begin{exercise}[$\star\star$]\label{exo:g7:triangles:9}
Dua sisi sebuah segitiga berukuran $8$ cm dan $3$ cm. Di antara dua
nilai berapa sisi ketiganya harus terletak? Sebutkan satu panjang
bilangan bulat yang berhasil dan satu yang tidak.
\end{exercise}

\begin{exercise}[$\star\star$]\label{exo:g7:triangles:10}
Pada segitiga $ABC$, $\widehat A = 2\widehat B$ dan
$\widehat C = 90^\circ$. Hitunglah $\widehat A$ dan $\widehat B$.
\end{exercise}

\begin{exercise}[$\star\star\star$]\label{exo:g7:triangles:11}
Sudut \omterm{def:g4:geometry:polygon}{segi empat} $ABCD$ mana pun dapat dipelajari dengan memotongnya
sepanjang sebuah diagonal menjadi dua segitiga. Pakailah ini untuk
membuktikan bahwa keempat sudut setiap \omterm{def:g4:geometry:polygon}{segi empat} berjumlah
$360^\circ$, lalu simpulkan \omterm{def:g1:addition:def}{jumlah} sudut \omterm{def:g4:geometry:polygon}{segi lima} (lima sisi).
\end{exercise}

\section{Soal: Berjalan mengelilingi blok}

\begin{problem}[{Soal akhir pekan --- sudut segi banyak mana pun:
$(n-2) \times 180^\circ$ di dalam, selalu $360^\circ$ putaran di
luar, dan mengapa hanya tiga bangun beraturan yang mengubin lantai}]
\label{pb:g7:triangles:1}
\Cref{exo:g7:triangles:11} memotong \omterm{def:g4:geometry:polygon}{segi empat} menjadi dua segitiga
lalu menemukan $360^\circ$. Soal ini mendorong gagasannya ke \omterm{def:g4:geometry:polygon}{segi
banyak} dengan sisi sebanyak apa pun, lalu menemukan bukti kedua yang
sama sekali berbeda --- dengan \emph{berjalan} mengelilingi bangunnya
dan menjumlahkan belokannya --- dan berakhir di lantai kamar mandimu:
di antara semua \omterm{def:g4:geometry:polygon}{segi banyak} beraturan, tepat tiga yang dapat mengubin
bidang tanpa celah atau tumpang tindih, dan lebah sudah memilih
miliknya.

\medskip
\noindent\textbf{Bagian I --- Sudut di dalam.}
\begin{enumerate}
  \item Dari satu \omterm{def:g5:solids:def}{titik sudut} \omterm{def:g4:geometry:polygon}{segi enam} (enam sisi), gambarlah semua
        diagonal yang berangkat dari \omterm{def:g5:solids:def}{titik sudut} itu. Menjadi berapa
        segitiga \omterm{def:g4:geometry:polygon}{segi enamnya} terpotong, dan berapa \omterm{def:g1:addition:def}{jumlah}
        \emph{semua} sudut dalamnya
        (\cref{thm:g7:triangles:sum})?
  \item Perumumlah: dari satu \omterm{def:g5:solids:def}{titik sudut} \omterm{def:g4:geometry:polygon}{segi banyak} dengan $n$ sisi,
        berapa banyak segitiga yang dihasilkan kipas diagonalnya?
        Simpulkan rumus \omterm{def:g1:addition:def}{jumlah} sudut dalamnya, lalu periksalah untuk
        $n = 3$, $4$, $5$, $6$.
  \item Hitunglah \omterm{def:g1:addition:def}{jumlah} sudut dalam sebuah segi sepuluh
        ($10$ sisi) dan sebuah segi dua belas ($12$ sisi).
  \item Pada \omterm{def:g4:geometry:polygon}{segi banyak} \emph{beraturan} semua sudut dalamnya sama
        besar. Hitunglah sudut dalam \omterm{def:g6:shapes:triangles}{segitiga sama sisi}, persegi, \omterm{def:g4:geometry:polygon}{segi
        lima} beraturan, \omterm{def:g4:geometry:polygon}{segi enam} beraturan dan \omterm{def:g4:geometry:polygon}{segi delapan}
        beraturan.
  \item Jelaskan, tanpa rumus apa pun, mengapa \omterm{def:g1:addition:def}{jumlah} sudut itu
        bertambah tepat $180^\circ$ setiap kali \omterm{def:g4:geometry:polygon}{segi banyaknya}
        bertambah satu \omterm{def:g5:solids:def}{titik sudut}.
\end{enumerate}

\medskip
\noindent\textbf{Bagian II --- Berjalan mengelilingi blok.}
Berjalanlah sepanjang tepi sebuah blok bersegi banyak, selalu maju. Di
setiap pojoknya kamu \emph{berbelok} sejauh suatu sudut --- \emph{sudut
luar} pojok itu; belokannya berlanjut sampai kamu kembali ke titik
awalmu, menghadap ke arah awalmu.
\begin{enumerate}[resume]
  \item Pada pojok yang sudut dalamnya $\widehat A$, sejauh berapa
        kamu berbelok? (Sudut dalam dan sudut luarnya berpelurus.)
        Hitunglah ketiga belokan segitiga dengan sudut $40^\circ$,
        $60^\circ$, $80^\circ$.
  \item Jumlahkan ketiga belokan pertanyaan 6. Kenyataan tentang satu
        lingkaran penuh apa yang kamu amati?
  \item Jelaskan mengapa belokan sepanjang perjalanan lengkap
        mengelilingi \omterm{def:g4:geometry:polygon}{segi banyak} \emph{mana pun} --- tiga sisi atau
        tiga puluh --- selalu berjumlah tepat $360^\circ$. (Apa yang
        sudah dilakukan hidungmu, dengan seluruh belokan digabungkan,
        ketika kamu tiba kembali?)
  \item Simpulkan rumus sudut dalamnya untuk kedua kalinya: pada
        setiap dari $n$ pojoknya, sudut dalam $+$ belokan
        $= 180^\circ$; jumlahkan atas semua pojoknya lalu pakailah
        pertanyaan 8.
  \item Pada \omterm{def:g4:geometry:polygon}{segi banyak} beraturan semua belokannya sama besar, jadi
        setiap sudut luarnya $\frac{360^\circ}{n}$. Pakailah ini untuk
        menjawab seketika: \omterm{def:g4:geometry:polygon}{segi banyak} beraturan mana yang sudut
        dalamnya $150^\circ$? Dan mengapa \emph{tidak ada} \omterm{def:g4:geometry:polygon}{segi banyak}
        beraturan yang sudut dalamnya $155^\circ$?
\end{enumerate}

\medskip
\noindent\textbf{Bagian III --- Mengubin lantai.}
Di sekitar setiap titik pada lantai berubin, pojok ubin yang bertemu
harus berjumlah tepat $360^\circ$ --- tanpa celah, tanpa tumpang
tindih (\cref{pb:g6:lines:1}).
\begin{enumerate}[resume]
  \item Untuk masing-masing dari \omterm{def:g6:shapes:triangles}{segitiga sama sisi}, persegi dan \omterm{def:g4:geometry:polygon}{segi
        enam} beraturan: berapa banyak salinan yang bertemu di satu
        titik pojok pengubinannya? Periksalah $360^\circ$-nya setiap
        kali.
  \item Tunjukkan bahwa \omterm{def:g4:geometry:polygon}{segi lima} beraturan tidak dapat mengubin
        lantai: berapa \omterm{def:g1:addition:def}{jumlah} tiga pojoknya, dan berapa \omterm{def:g1:addition:def}{jumlah} empat
        pojoknya?
  \item Tunjukkan bahwa tidak ada \omterm{def:g4:geometry:polygon}{segi banyak} beraturan dengan
        \emph{lebih} dari enam sisi yang dapat mengubin: sedikitnya
        tiga ubin harus bertemu di setiap pojok (setiap sudutnya
        kurang dari $180^\circ$), lalu apa yang terjadi pada tiga
        pojok berukuran $135^\circ$ atau lebih? Simpulkan daftar
        lengkap pengubin beraturannya.
  \item Lantai kamar mandi sering mencampur \omterm{def:g4:geometry:polygon}{segi delapan} yang
        beraturan dengan persegi kecil. Periksalah bahwa pojok ini berhasil:
        $135^\circ + 135^\circ + 90^\circ$. Campuran klasik yang lain
        menggabungkan persegi, \omterm{def:g4:geometry:polygon}{segi enam} beraturan dan segi dua belas
        beraturan di setiap pojoknya: periksalah itu juga (sudut dalam
        segi dua belasnya menyusul dari cara pertanyaan 10).
  \item Pilihan lebah: sel sarang lebah berbentuk \omterm{def:g4:geometry:polygon}{segi enam} beraturan.
        Periksalah pojoknya ($3$ sudut), lalu jelaskan dalam satu
        kalimat --- dengan bantuan penemuan Dido
        (\cref{pb:g6:measure:1}) --- mengapa, dari ketiga ubin
        beraturan yang mungkin, \omterm{def:g4:geometry:polygon}{segi enamnya} yang paling murah
        lilinnya untuk madu yang ditampungnya.

\end{enumerate}
\end{problem}
